\section{Background and Problem Setting}
\label{sec:preliminaries}
\subsection{PWM-Driven Vehicular Actuation}
\begin{itemize}
    \item Canonical architecture: controller ECU $\rightarrow$ PWM peripheral $\rightarrow$ interconnect/isolation $\rightarrow$ gate driver $\rightarrow$ power stage $\rightarrow$ plant $\rightarrow$ sensors.
    \item Timing domains: control-loop period $T_c$ vs PWM switching frequency $f_s$; why single-period perturbations can matter for safety.
    \item Observability options: PWM boundary signals, gate output/switch-node timing features, and plant-side switching signatures (e.g., current ripple).
\end{itemize}

\subsection{Threat Landscape: Actuation-Path Attack Classes}
\begin{itemize}
    \item \textbf{Actuation-path tampering} (post-compute, pre-switch): pulse insert/drop, edge skew, effective dead-time distortion; includes EMI/false-actuation-style attacks.
    \item \textbf{PWM peripheral/register corruption}: DMA-like overwrite, rogue interrupt/task, or transient faults causing incorrect compare/dead-time values despite correct control computation.
    \item Positioning: these classes are motivated by prior academic/real-world demonstrations; our work targets detection/attestation of \emph{realized actuation}, not network command authentication.
\end{itemize}

\subsection{Problem Statement and Security Goal}
\begin{itemize}
    \item Let the controller compute control intent $u_k$ and PWM parameters $d_k$ at step $k$ using measured state $x_k$.
    \item The system emits a digital PWM waveform $p(t)$, which is transformed by the actuation path into physical switching behavior and plant response.
    \item \textbf{Actuation Integrity (goal):} detect when realized actuation deviates from intended $d_k$ beyond tolerance $\epsilon$, within bounded detection latency $\Delta$ (e.g., within a small number of control cycles), while controlling false positives/negatives.
\end{itemize}

\subsection{System Model}
\begin{itemize}
    \item Components: (i) motor-control ECU (trusted control task), (ii) actuation path (potentially untrusted), (iii) plant + sensors, and (iv) separate verifier/safety supervisor ECU (trusted, remote-from-controller).
    \item Signals: control intent $u_k$, state $x_k$, intended switching parameters $d_k$, evidence $E_k$ (measured at a boundary or via switching-coupled plant signatures).
    \item Verifier access: authenticated witness of intended command (at minimum $d_k$, optionally $(u_k,x_k)$) and independent evidence $E_k$.
\end{itemize}

\subsection{Threat Model}
\begin{itemize}
    \item Adversary objective: induce unsafe behavior by causing realized actuation to differ from intended $d_k$, including transient (single/few-cycle) manipulations.
    \item Adversary capabilities (primary):
    \begin{enumerate}
        \item \emph{Actuation-path tampering} between PWM generation and switching (wireless EMI/FAI or conducted superposition).
        \item \emph{PWM peripheral/register corruption} that alters emitted PWM without changing the control algorithm.
    \end{enumerate}
    \item Trusted computing base (TCB): verifier/safety supervisor, evidence acquisition path, and authenticity of the witness export of $d_k$ (and optional state summaries).
    \item Out of scope: attacker compromise of verifier/evidence acquisition; full arbitrary code execution on the controller that can suppress alarms or forge reports.
\end{itemize}